\chapter{Check each representation of the PNG}
\MARSChapterHeading{Check each representation of the PNG}
\label{revision-lab:png}

Chapter 24 follows a complete decoder. This lab asks you to establish what each
stage proves, then implement a forward row filter from original RGB samples.
Import the native lesson \emph{Check each PNG representation}, built from
\path{revision/lessons/033-png}. Its six complete worked experiments lead to one
continuing source, \texttt{independent/png-filter-scanline.asm}. Only that
source's forward-filter entry is unfinished.

Begin with \texttt{worked/pipeline-boundaries.asm}. It embeds the exact
115-byte trace and retains the supplied decoder functions and format tables.
Its smaller caller uses static buffers: \texttt{lesson\_stream},
\texttt{lesson\_rows} and \texttt{lesson\_frame}. Required observations need
no host input or path edit. The framebuffer stride is seventeen words and its
selected extent is 280 bytes. The separate \texttt{worked/packing-tail.asm}
experiment crosses the sixteen-pixel vector boundary.

Answer predictions inline, then use ordinary Debug at the named boundaries.
Inspect state before executing the anchored instruction. Finish the process
before starting the next activity. The lesson's six Debug activities contain
eighteen stops, including full-width MSA input and packed-output values.
After exploration with edits or backsteps, start a fresh forward Debug session
to record the required sequence. Hints and free-writing availability follow
your instructor's release policy.

The standalone worked caller remains at
\path{revision/labs/files/png-viewer.asm}, with
\texttt{png-filter-scanline.asm} beside it for independent work. That viewer
loads a host file and allocates buffers on the heap. The native lesson's smaller
caller and the viewer use the same decoder helpers, but their buffer-location
variables differ. Use the workflow named by each observation below.

For the larger capstone assignment, the original
\texttt{packages/033-png/starter/main.asm} leaves the loader, collector, inflater,
scalar/MSA reconstruction, packing and complete decode entries unfinished.
The construction directory retains staged attachments and observations for
those entries. The bit-reader and table-construction work below extends that
capstone; the independent forward filter does not require implementing an inflater. Follow the published interfaces; private register choices,
frames and equivalent algorithms may differ. The worked program is available
for inspection, but the independent task below requires you to reason in the
opposite transformation direction.

\section{Make a boundary ledger before debugging}

Use the embedded two-by-five trace in the native lesson, or generate it for
the standalone viewer. Record file extent, each chunk's header and
payload extent, joined stream count, expected row count and active pixel count.
Your ledger should explain 115, 46, 35, 30 and 40 without calling them all the
image size. Calculate the compressed bytes contributed by each IDAT. Explain
why concatenation must precede zlib interpretation when the header is split
between the two chunks.

Follow the native caller at \texttt{collector\_return},
\texttt{inflate\_return}, \texttt{publication\_rows\_return} and
\texttt{publication\_blit\_return}. Their successful \texttt{v1} counts are
46, 35, 35 and 40. Reconstruction retains the five selectors, so it returns
the 35-byte row-buffer extent even though that buffer contains thirty RGB
sample bytes. At \texttt{complete\_image}, \texttt{s0}, \texttt{s1} and
\texttt{s2} contain the published width two, height five and forty active bytes.
The second row starts 68 bytes after the first; \texttt{s5} reads the unchanged
first byte of the first padding gap. One sentinel does not prove the complete
access contract, which the independent grading observes separately.

In the standalone viewer, at \texttt{collector\_return} the decoder has a private successful
collector descriptor. Its capacity is not a width field. The public complete
decode context has not yet published its result fields. Inspect
\texttt{compressed\_address}, then follow that stored pointer to the joined
bytes. Distinguish the address of this caller variable from the address stored
inside it. Do the same for \texttt{row\_address} and
\texttt{framebuffer\_address} rather than assuming one fixed heap address for
all configurations.

Change a payload byte without updating its chunk CRC. Predict the stage that
will fail and which public objects must remain unchanged. Then repair the CRC
but make the compressed stream invalid. The collector may now succeed while
inflation fails. Those experiments test different promises; neither should
publish pixels or complete-decoder context results.

\section{Build a bit reader and a table decoder deliberately}

In the native lesson, \texttt{worked/bit-interpretation.asm} isolates the
arriving bits one, zero. They give numerical field value one but developed
canonical-code value binary \texttt{10}. Its second pause decodes stored length
35 and complement 65500, whose sixteen-bit XOR is hexadecimal \texttt{ffff}.
Use these small observations to explain the interpretation before attempting
the larger table decoder.

For the stored trace, identify the final flag, block type, byte alignment,
length and complement. Test a bad complement, a truncated payload and a valid
empty stored block separately. The general inflater can accept expected output
zero even though a supported PNG's positive dimensions predict nonzero rows.
Require the exact produced count before accepting its trailer.

For compressed blocks, derive canonical codes from code lengths. Count symbols
of each length and order equal-length symbols by numeric value. Track available
prefix slots with $K_0=1$ and
\[
 K_\ell=2K_{\ell-1}-count_\ell.
\]
A negative value means too many codes. A complete tree finishes at zero.
Remaining slots require one of the explicitly permitted sparse cases: a single
one-bit literal/length or distance code, or an empty distance alphabet that
cannot serve a match. The code-length alphabet must be complete. An unused
branch is never a valid decoded symbol, and the literal/length alphabet must
include end-of-block symbol 256.

Use the main chapter's three-symbol alphabet to decode A, length three,
distance one, end-of-block. Write the arriving bits and the produced bytes.
Then compare code \texttt{10} with the numerical field extracted from the same
two arriving bits. That small discrepancy should explain which helper can read
extra-bit counts and which algorithm must decode symbols.

Dynamic headers expand repeat operations. Symbol 16 repeats the previous length
three through six times; 17 supplies three through ten zeros; 18 supplies eleven
through 138 zeros. Their extra bits select a count within those ranges. Before
expanding, compare the selected count with remaining array positions. Test a
repeat without a previous length and a repeat that would exceed the final
position by one. Reset per-block tables and retain only the produced-byte
history across blocks. A successful fixed block followed by a malformed dynamic
block must not reuse a stale table to make the latter appear valid.

The original construction material supplies the base/extra-bit tables and
separate inflater cases. Keep those syntax tables at hand while implementing;
they need not be memorized. Test reserved literal/length and distance symbols
before using them as table indices.

\section{Trace history and row dependencies}

Debug \texttt{worked/overlapping-history.asm} in the native lesson. It begins
with one produced byte A and requests six new bytes at distance one.
Write the output and source cursor after each store. Explain why loading six
source bytes in advance would consume unproduced data. Then test distance
larger than produced history and length larger than remaining output. A mapped
address is not proof of valid history or selected output capacity.

Pause the complete program at \texttt{inflate\_return}. Interpret the buffer
as five groups of seven bytes. For the first Up row, derive samples 15, 25, 35,
45, 55 and 65. Continue to \texttt{publication\_rows\_return} and compare.
The selectors remain in place; sample values change. For the Sub row, follow
one red channel across pixels and identify the reconstructed left neighbor
before using it.

The native \texttt{worked/predictor-precision.asm} keeps both the wide Average
and a prematurely byte-wrapped calculation visible. For unsigned neighbors
255 and 145, the sum is 400 and the predictor is 200; wrapping first would
produce 72. A current sample 201 therefore emits residual one. The final
residual wraps modulo 256, while predictor intermediates keep their precision.
For Paeth, test $a=255,b=0,c=85$. The estimate is 170, and $a$ ties $c$ at
distance 85; $a$ wins. For $a=0,b=3,c=1$, $b$ ties $c$ at distance one and
wins. Changing a strict comparison in a tie decision can corrupt only a few
pixels, so a typical image is a weak test of that policy.

Make the final selector five while keeping outer and Adler checksums valid.
The reconstruction function must detect every selector before changing any
samples. The complete decoder must leave framebuffer and context unchanged.
Neither guarantee requires restoring the inflater's scratch bytes to their
old contents. Describe failure boundaries by object, rather than using an
unqualified promise that all memory rolls back.

\section{Explain the vector result at both widths}

Use \texttt{worked/packing-tail.asm} in the native lesson. At
\texttt{pack\_inputs\_loaded}, inspect all 128 bits of \texttt{w0} through
\texttt{w2}: their byte lanes contain samples zero through 47. Lane zero is
at the low end of the displayed value. At \texttt{pack\_first\_written},
\texttt{w3} contains four numerical words for pixels zero through three.
At \texttt{pack\_last\_written}, it contains pixels twelve through fifteen.
The group cursor has not yet advanced at either pause. The final
\texttt{packed\_tail} observation checks pixels five, fifteen and sixteen,
68 active bytes and the first retained guard byte.

Figure~\ref{fig:revision-png-native-vector} shows the first packed group in the
actual debugger. Use the lane interpretation to check each pixel separately.

\begin{figure}[!ht]
  \centering
  \begin{tikzpicture}[alt={Three crops of one actual light-theme debugger photograph. A identifies the paused frame at packing-tail.asm line 1607, before the next no-operation instruction. B shows the complete written w3 vector. C interprets it as four 32-bit integer lanes: hexadecimal 00000102, 00030405, 00060708 and 00090a0b, starting with lane zero. These are the first four numerical RGB pixel words.}]
    \node[anchor=north west,inner sep=0,font=\small\bfseries] (title) at (0,0) {Read four packed pixels before advancing the group};
    \node[anchor=north west,inner sep=0] (frame) at ([yshift=-3mm]title.south west) {\MARSRevisionCrop{revision/captures/png/png-native-vector.png}{0.376\linewidth}{0.03828125}{0.22222222}{0.77812500}{0.75111111}};
    \node[anchor=east,font=\small\bfseries] at ([xshift=-1mm]frame.west) {A};
    \node[anchor=north west,inner sep=0] (vector) at ([yshift=-3mm]frame.south west) {\MARSRevisionCrop{revision/captures/png/png-native-vector.png}{0.96\linewidth}{0.24843750}{0.09777778}{0.28281250}{0.87777778}};
    \node[anchor=east,font=\small\bfseries] at ([xshift=-1mm]vector.west) {B};
    \node[anchor=north west,inner sep=0] (lanes) at ([yshift=-3mm]vector.south west) {\MARSRevisionCrop{revision/captures/png/png-native-vector.png}{0.96\linewidth}{0.26171875}{0.03333333}{0.26953125}{0.94333333}};
    \node[anchor=east,font=\small\bfseries] at ([xshift=-1mm]lanes.west) {C};
  \end{tikzpicture}
  \caption{The debugger is paused at \texttt{pack\_first\_written}, line 1607 (A).
    The preceding store has written sixteen output bytes, and \texttt{w3} retains
    four packed pixels (B). Read the 32-bit lanes from lane zero (C):
    \texttt{0x00000102} means red zero, green one and blue two;
    \texttt{0x00030405} means red three, green four and blue five.
    The high byte of each numerical word is zero. A vector's printed hexadecimal
    value places its low lane on the right, while the lane list begins with lane
    zero. All three crops are from the same unchanged photograph.}
  \label{fig:revision-png-native-vector}
\end{figure}


Then compare widths fifteen, sixteen and seventeen. For packing, mark exactly which
48 RGB bytes feed a complete sixteen-pixel group, which four output vectors
receive them, and which samples the tail handles. Derive the mask for pixel
five across input vectors zero and one. Explain why the shuffle destination
must be restored from a retained mask before reuse.

Compare numerical framebuffer words under both byte orders. The first trace
pixel remains \texttt{0x000a141e} even though its four memory bytes differ.
With stride seventeen, mark the five eight-byte active rows and the four
sixty-byte gaps. Derive the required extent 280, and explain why 340 would
incorrectly require a final padding gap. Check those gaps for retained bytes;
a correct active pixel count does not authorize stores into padding.

Up reconstruction can add independent byte groups from the completed previous
row. Sub reconstruction depends on new values three bytes earlier in its
current row. State those dependencies before proposing vector code. When
comparing scalar and MSA modes, use identical file bytes and stride, establish
exact output equality, then separate complete decode work from loading and
assembly. The original observation tables retain their historical measurement
scope; reproduce a measurement before treating it as a result for your program.

\section{Implement the forward row filter}

Replace the intentional \texttt{break} at \texttt{png\_filter\_scanline} in
the independent file. Arguments are current RGB pointer in \texttt{a0}, previous
RGB pointer in \texttt{a1}, output pointer in \texttt{a2}, and settings pointer
in \texttt{a3}. Settings are three aligned native unsigned doublewords: width,
filter selector and output capacity. Width must be one through 512 and selector
zero through four. The required output length is $1+3w$. Capacity is unsigned;
a high-bit-set value is large, rather than a negative shortage.

After validating settings, emit the selector followed by $3w$ residual bytes.
Use the original current-row sample three bytes earlier as $a$, the previous-row
sample as $b$, and its preceding same-channel sample as $c$. Missing left
neighbors are zero. A null previous pointer means the first image row, so $b$
and $c$ are zero. Filters None and Sub never read the previous pointer, even
if it is nonnull. Preserve full predictor precision and the stated Paeth tie
order before subtracting modulo 256.

The RGB byte intervals may be unaligned and may coincide with each other,
since both are read-only. Output must be disjoint from both inputs and settings;
settings must be disjoint from the RGB intervals. Only the first 24 settings
bytes and the selected input/output intervals belong to the call. Return zero
and $1+3w$ on success. An invalid setting returns one and zero before accessing
any RGB or output byte and changes no caller object. Preserve every bit of
\texttt{s0} through \texttt{s7}, \texttt{fp}, \texttt{sp} and the return
continuation. The incoming stack is eight-byte aligned. Do no I/O in the entry.

Before implementation, Debug \texttt{worked/native-settings.asm}.
Inspect the full width, hexadecimal \texttt{0000000100000001}.
The low word alone looks like one, but the complete value is outside the accepted range. The second
pause compares unsigned capacity hexadecimal \texttt{8000000000000000}
with seven. The unsigned shortage test is false. Native assessments assemble
your source and settings in both byte orders, so native data declarations must
retain their intended numerical meaning in either representation.

The supplied demo uses width two, Average, current RGB
\texttt{25,35,45,55,65,75} and previous RGB \texttt{20,30,40,50,60,70}.
It must produce seven hexadecimal bytes
\texttt{03 0f 14 19 12 12 12}. Pause at \texttt{filter\_demo\_return} to see
\texttt{v0/v1} before the caller records them in \texttt{filter\_status} and
\texttt{filter\_count}. The expected count is seven bytes, rather than six
samples or two pixels.

Start with None, then Sub and Up, then Average and Paeth. Compare a forward
loop with a backward loop: both can read original neighbors because neither
writes its input. This is a different dependency proof from in-place unfiltering.
Do not add exact-alias output support silently; that would require preserving
original neighbors before output stores can overwrite them.

\section{Check exact bytes and the inverse}

Use \emph{Check this program} in the native lesson. Each predictor stage
retains 44 original cases in both byte orders; the boundary stage retains
46 more, including six calls sharing one machine. The final assessment rechecks
all 266 on the same source. When edits make earlier evidence stale, follow the
renewal prompts. Prediction attempts and revealed hints remain saved. No exact
loop, scratch-register choice or private frame layout is required.

For standalone CLI work, run the public independent feedback from the book repository:
\Needspace{8\baselineskip}
\MARSCodeLabel{Grade the new forward transformation}
\begin{Verbatim}
java -jar "$CLI" grade \
  revision/labs/checks/png-filter-scanline.kts \
  revision/labs/files/png-filter-scanline.asm
\end{Verbatim}
The check covers every filter, top-row and previous-row cases, both byte
orders, unaligned byte intervals, boundary widths, full-width invalid settings,
unsigned capacities, coincident read-only inputs, repeated calls and explicit
Paeth ties. It compares exact residuals and retained caller state, then applies
an inverse recurrence to recover every original sample. A matching checksum
alone would be weaker evidence.

For the original complete decoder, keep its separate harness out of your source
and use its supplied runner:
\Needspace{9\baselineskip}
\MARSCodeLabel{Check the complete pipeline with separate grading storage}
\begin{Verbatim}
python3 packages/033-png/grade.py \
  revision/labs/files/png-viewer.asm --domain decode \
  --cli "$CLI" --java "$JAVA"
\end{Verbatim}
Here \texttt{JAVA} is the full JDK 25 Java executable. Omit \texttt{--domain}
to check all eight domains. The runner supplies \texttt{checks/harness.asm}
as another assembly source. Preserve the public callable entries and reserve
the \texttt{pngCheck\_} prefix for those resources. The independent forward
function has its own feedback; the original decoder runner does not grade it.

Finally collect five forward-filtered rows in order, using a null previous
pointer for the first and the preceding original RGB row for each later one.
Export those 35 binary bytes as \texttt{rows.bin}; include each selector, with
no text addresses, spaces or hexadecimal characters. The supplied host wrapper
provides compression and container writing without implementing your filter:
\Needspace{8\baselineskip}
\MARSCodeLabel{Wrap the simulator-produced rows in a complete PNG}
\begin{Verbatim}
python3 revision/labs/files/png-from-rows.py rows.bin result.png \
  --width 2 --height 5 --compression fixed
\end{Verbatim}
The imported lesson retains the same tools under \texttt{tools/}; run those
copies from its folder, using fresh output paths. The command above uses the
book repository's path. The wrapper validates dimensions, exact row count and selectors, then creates one zlib
stream split across two consecutive IDAT chunks. It refuses an existing output.
Keep \texttt{result.png} separate from your original input, then select it for
the decoder or place a copy named \texttt{image.png} in a fresh working directory.
Decode with the worked program and an independent image reader and compare
exact RGB samples. This separates your prediction logic from compression and
gives an artifact whose pixels can be checked outside MARS.
